% Part: incompleteness
% Chapter: incompleteness-provability
% Section: godels-paper

\documentclass[../../../include/open-logic-section]{subfiles}

\begin{document}

\olfileid{inc}{inp}{gop}

\olsection{د ګوډل له اصلي مقالې سره پرتله}

د ګوډل د ۱۹۳۱ز کال مقالې ته لږ وخت ورکول ګټور دي۔ د هغې پېژندنه
د هغو مفکورو لنډ بيان کوي چې موږ همدا اوس وڅېړلې۔ که مقاله يوازې
په تېره هم وګورئ، بيا هم د هر پړاو کار څرګندېږي: لومړے ګوډل صوري
نظام~$P$ بيانوي (نحو، بديهي اصول او د ثبوت قاعدې)؛ بيا ابتدايي
بازګشتي تابعې او اړيکې تعريفوي؛ ورپسې ښيي چې~$x B y$ ابتدايي
بازګشتي ده او استدلال کوي چې ابتدايي بازګشتي تابعې او اړيکې
په~$\Th{P}$ کښې تمثيلېږي۔ بيا د ناتکميلۍ قضيه لکه پورته ثابتوي۔
په درېيمه برخه کښې ښيي چې ناثابتېدونکې دعوه د حساب د ژبې يوه جمله
اخيستے شو۔ د~$\beta$-لمې اصل همدلته دے؛ موږ هم دا لمه د لړيو
د سمبالولو دپاره وکاروله، کله چې مو ښودله چې بازګشتي تابعې
په~$\Th{Q}$ کښې تمثيلېږي۔ ګوډل دومره وړاندې نۀ ځي چې د بسيا
بديهي اصولو يو لږترين سټ بېل کړي، خو اوس پوهېږو چې~$\Th{Q}$ د
دې کار دپاره بسيا ده۔ په پاې کښې، په څلورمه برخه کښې د ناتکميلۍ
د دويمې قضيې د ثبوت لنډه نقشه ورکوي۔

\end{document}
